\documentclass[a4paper,12pt]{article}

\usepackage[T2A]{fontenc}
\usepackage[utf8]{inputenc}
\usepackage[russian]{babel}

\usepackage{amsmath, amssymb, amsthm, mathtools}
\usepackage{helvet}
\renewcommand{\familydefault}{\sfdefault}
\usepackage{icomma}
\usepackage{geometry}
\usepackage{microtype}
\usepackage{tikz}
\usepackage{tkz-euclide}
\usepackage{pgfplots}
\pgfplotsset{compat=1.18}
\usepackage{tabularx, booktabs, longtable}
\usepackage{enumitem}
\usepackage{hyperref}
\usepackage{parskip}
\usepackage{fancyhdr}
\usepackage{setspace}
\usepackage{xcolor}

\definecolor{accent}{HTML}{008080}
\definecolor{darkaccent}{HTML}{004D4D}
\definecolor{textgray}{HTML}{333333}

\geometry{left=2cm, right=2cm, top=2cm, bottom=2cm}
\onehalfspacing

\hypersetup{
    colorlinks=true,
    linkcolor=accent,
    filecolor=magenta,
    urlcolor=cyan,
}

\pagestyle{fancy}
\fancyhf{}
\rhead{\thepage}
\renewcommand{\headrulewidth}{0.5pt}

\begin{document}

% Титульный лист
\begin{titlepage}
    \centering
    \vspace*{2cm}

    {\Huge\bfseries\color{accent} Чек-лист}\par
    \vspace{0.5cm}
    {\Large\color{textgray} Аннуитетная схема платежей: экономические задачи}\par

    \vspace{3cm}

    {\large\color{textgray} Лёвин Артём Александрович}\par
    {\large\color{textgray} эксклюзивно для Макса}\par

    \vspace{2cm}

    {\large\color{textgray} \today}\par

    \vfill

    % Кривая Безье внизу титульного листа
    \begin{tikzpicture}
        \draw[line width=2pt, color=accent, opacity=0.6]
            (0,0) .. controls (2,1.5) and (4,-1) .. (6,0.5)
            .. controls (8,2) and (10,-0.5) .. (12,0);
    \end{tikzpicture}

\end{titlepage}

\newpage

% Основной чек-лист
\section*{Ключевые знания по теме}

\subsection*{Основные понятия и определения}
\begin{itemize}[label=\textcolor{accent}{\textbullet}, leftmargin=*]
    \item \textbf{Аннуитетная схема} — схема кредитования, при которой выплачиваются \textit{равные платежи} через равные промежутки времени
    \item \textbf{Начальная сумма кредита} ($S$) — сумма, взятая в кредит
    \item \textbf{Процентная ставка} ($p\%$) — процент, под который берётся кредит; в расчётах используется $r = 1 + \frac{p}{100}$
    \item \textbf{Платёж} ($x$) — одинаковая сумма, выплачиваемая каждый период
    \item \textbf{Остаток задолженности} — сумма, остающаяся должной после очередного платежа
    \item \textbf{Переплата} — разница между общей суммой выплат и начальной суммой кредита
    \item \textbf{Процент переплаты} — отношение переплаты к начальной сумме, выраженное в процентах
\end{itemize}

\subsection*{Математическая модель (таблица)}
Для аннуитетной схемы составляется таблица следующего вида:

\begin{center}
\renewcommand{\arraystretch}{1.4}
\begin{tabular}{ccccc}
\toprule
Год & Начальная сумма & Начисленный \% & Общая задолженность & Платёж & Остаток \\
\midrule
1 & $S$ & $Sp$ & $S + Sp = Sr$ & $x$ & $Sr - x$ \\
2 & $Sr - x$ & $(Sr-x)p$ & $(Sr-x)r$ & $x$ & $(Sr-x)r - x$ \\
3 & $(Sr-x)r - x$ & $\ldots$ & $\ldots$ & $x$ & $\ldots$ \\
\bottomrule
\end{tabular}
\end{center}

где $r = 1 + \frac{p}{100}$ — коэффициент увеличения задолженности.

\subsection*{Вывод формулы для платежа}
Из условия, что после последнего платежа остаток равен нулю, получаем уравнение. Для $n=3$ лет:

$$Sr^3 - xr^2 - xr - x = 0$$

Отсюда выражаем величину одного платежа:

$$x = \frac{Sr^3}{r^2 + r + 1}$$

В общем случае для $n$ лет:

$$x = \frac{Sr^n}{r^{n-1} + r^{n-2} + \ldots + r + 1} = \frac{Sr^n(r-1)}{r^n - 1}$$

\subsection*{Формулы для расчётов}
\begin{itemize}[label=\textcolor{accent}{\textbullet}, leftmargin=*]
    \item \textbf{Общая сумма выплат:} $X_{общ} = n \cdot x$
    \item \textbf{Переплата:} $P = X_{общ} - S = nx - S$
    \item \textbf{Процент переплаты:} 
    $$\frac{P}{S} \cdot 100\% = \left(\frac{nr^n(r-1)}{r^n-1} - 1\right) \cdot 100\%$$
    \item Важно: процент переплаты \textit{не зависит} от начальной суммы $S$, а зависит только от $n$ и $p$
\end{itemize}

\subsection*{Методы решения задач}
\begin{enumerate}[label=\arabic*., leftmargin=*]
    \item Составить таблицу для каждого года кредитования
    \item Ввести обозначения: $S$, $p\%$, $r = 1 + \frac{p}{100}$, $x$ — платёж
    \item Записать рекуррентную формулу для остатка задолженности
    \item Приравнять остаток после последнего платежа к нулю
    \item Выразить $x$ из полученного уравнения
    \item При работе с двумя разными сроками кредитования — составить две таблицы и поделить формулы друг на друга для исключения $S$
    \item Использовать замену переменной при решении уравнений высоких степеней
    \item Большие числа раскладывать на множители для упрощения вычислений
\end{enumerate}

\subsection*{Типичные ошибки}
\begin{itemize}[label=\textcolor{accent}{\textbullet}, leftmargin=*]
    \item Забывают перевести проценты в десятичную дробь: $p\% \to \frac{p}{100}$
    \item Не делают замену $r = 1 + \frac{p}{100}$, что приводит к громоздким выкладкам
    \item Ошибаются при раскрытии скобок в рекуррентных формулах
    \item Лишний раз умножают на $x$ при составлении уравнения
    \item Не сокращают $S$ при делении формул друг на друга
    \item Не проверяют здравый смысл полученных чисел
    \item Пытаются считать большие числа без предварительного разложения на множители
\end{itemize}

\subsection*{Примеры решения}

\textbf{Пример 1.} Кредит 120 тыс. руб. под 20\% годовых на 3 года. Найти платёж.

\textit{Решение:} $S = 120$, $p = 20\%$, $r = 1{,}2$. 

$$x = \frac{120 \cdot 1{,}2^3}{1{,}2^2 + 1{,}2 + 1} = \frac{120 \cdot 1{,}728}{1{,}44 + 1{,}2 + 1} = \frac{207{,}36}{3{,}64} \approx 57$$

\textbf{Пример 2.} Сравнение кредитов на 2 и 4 года. Платежи относятся как $x_4 : x_2 = 2592 : 4392$. Найти $p$.

\textit{Решение:} Делим формулы:
$$\frac{x_4}{x_2} = \frac{r^4}{r^3 + r^2 + r + 1} \cdot \frac{r + 1}{r^2} = \frac{r^4}{(r^2+1)(r+1)} \cdot \frac{r+1}{r^2} = \frac{r^2}{r^2+1}$$

Получаем уравнение $\frac{r^2}{r^2+1} = \frac{2592}{4392} = \frac{36}{61}$, откуда $r^2 = \frac{36}{25}$, $r = 1{,}2$, $p = 20\%$.

\newpage

% Тренировочный блок
\section*{Тренировочный блок}

\subsection*{Базовый уровень (3 задания)}
\begin{enumerate}[label=\arabic*., leftmargin=*]
    \item Кредит 100 тыс. руб. взят под 10\% годовых на 2 года с аннуитетными платежами. Найдите величину ежегодного платежа.
    
    \item Найдите коэффициент $r$, если процентная ставка равна 15\%.
    
    \item Запишите формулу для остатка задолженности после первого платежа, если начальная сумма $S$, коэффициент $r$, платёж $x$.
\end{enumerate}

\subsection*{Средний уровень (3 задания)}
\begin{enumerate}[label=\arabic*., leftmargin=*, start=4]
    \item Кредит 200 тыс. руб. под 20\% годовых на 3 года. Найдите общую сумму выплат за весь срок кредитования.
    
    \item Докажите, что для $n=2$ лет формула платежа имеет вид: $x = \frac{Sr^2}{r+1}$.
    
    \item Найдите переплату по кредиту 150 тыс. руб. на 2 года под 10\% годовых с аннуитетными платежами.
\end{enumerate}

\subsection*{Уровень выше среднего (3 задания)}
\begin{enumerate}[label=\arabic*., leftmargin=*, start=7]
    \item Кредит на 3 года и кредит на 5 лет взяты под один и тот же процент. Отношение платежей $x_5 : x_3 = 3 : 2$. Найдите процентную ставку.
    
    \item Выведите формулу для процента переплаты при кредитовании на $n$ лет и покажите, что она не зависит от $S$.
    
    \item Кредит 500 тыс. руб. под 25\% годовых на 4 года. На сколько процентов переплата больше начальной суммы?
\end{enumerate}

\subsection*{Повышенный уровень (1 задание)}
\begin{enumerate}[label=\arabic*., leftmargin=*, start=10]
    \item Два кредита взяты под одну и ту же процентную ставку $p\%$. Первый кредит на 2 года с платежом 4 392 300 руб., второй на 4 года с платежом 2 592 300 руб. Найдите начальную сумму каждого кредита и процентную ставку.
\end{enumerate}

\newpage

% Ответы
\section*{Ответы к тренировочному блоку}

\begin{center}
\renewcommand{\arraystretch}{1.5}
\begin{tabular}{cl}
\toprule
\textbf{№ задания} & \textbf{Ответ} \\
\midrule
1 & $x \approx 57{,}62$ тыс. руб. \\
2 & $r = 1{,}15$ \\
3 & $Sr - x$ \\
4 & $X_{общ} \approx 322{,}4$ тыс. руб. \\
5 $Sr^2 = x(r+1) \Rightarrow x = \frac{Sr^2}{r+1}$ \\
6 & $P \approx 32{,}25$ тыс. руб. \\
7 & Уравнение: $\frac{r^2(r^2+1)}{r^4+r^3+r^2+r+1} = \frac{3}{2}$ \\
8 & $\frac{P}{S} = \frac{nr^n(r-1)}{r^n-1} - 1$ \\
9 & $\approx 144\%$ \\
10 & $p = 20\%$, $S_2 \approx 7{,}5$ млн руб., $S_4 \approx 7{,}5$ млн руб. \\
\bottomrule
\end{tabular}
\end{center}

\end{document}